\section*{Chimie (7 points)~: Étude de transformations chimiques}
En chimie des solutions, les transformations chimiques diffèrent selon les
couples intervenants et les conditions expérimentales. L'étude de ces
transformations qui correspondent dans certains cas à des réactions acide-base
ou des réactions d'oxydo-réduction peut se faire en utilisant des méthodes
différentes, ce qui permet de prévoir et de suivre l'évolution des systèmes
chimiques et de déterminer certaines grandeurs qui les caractérisent.

Cet exercice vise~:
\begin{itemize}
    \item l'étude cinétique d'une transformation chimique~;
    \item l'étude de transformations chimiques faisant intervenir des couples de types différents.
\end{itemize}

\subsection{Cinétique d'une transformation chimique}
On introduit dans un ballon une solution aqueuse $\mathrm{(S)}$ d'acide méthanoïque $\mathrm{HCO}_{2} \mathrm{H}$ de quantité de matière $n_{1}$. On y ajoute, à l'instant $t_{0}=0$, une solution aqueuse de dibrome $B r_{2(a q)}$ de quantité de matière $n_{2}$ (avec $n_{2}<n_{1}$ ). Le mélange réactionnel est maintenu à une température constante $\theta_{1}$. L'équation chimique de la réaction qui a eu lieu s'écrit~:

$$
\mathrm{HCO}_{2} \mathrm{H}_{(a q)}+\mathrm{Br}_{2(a q)} \rightarrow 2 \mathrm{Br}_{(a q)}^{-}+2 \mathrm{H}_{(a q)}^{+}+\mathrm{CO}_{2(\mathrm{~g})}
$$

Une étude expérimentale a permis de déterminer les valeurs de l'avancement $x$ et de la vitesse volumique $v$ de la réaction à différents instants $t$. Les résultats sont dressés dans le tableau cidessous.

\begin{center}
\begin{tabular}{|c|c|c|c|c|c|c|}
\hline
$t(\mathrm{~s})$ & 50 & 200 & 450 & 800 & 1231 & 1300 \\
\hline
$x\left(10^{-3} \mathrm{~mol}\right)$ & 0,190 & 0,600 & 0,941 & 1,13 & 1,20 & 1,20 \\
\hline
$v\left(10^{-6} \mathrm{~mol} . \mathrm{L}^{-1} . \mathrm{s}^{-1}\right)$ & 2,7 & 1,7 & 0,75 & 0,33 & 0 & 0 \\
\hline
\end{tabular}
\end{center}

1.1. En exploitant les données du tableau~:
a. déterminer la valeur de l'avancement final $x_{f}$ de la réaction. Déduire la valeur de $n_{2}$.
b. déterminer la valeur du temps de demi-réaction $t_{1 / 2}$.
c. interpréter qualitativement la variation de la vitesse volumique de la réaction.
1.2. On refait l'expérience en utilisant les mêmes quantités de matière à une température $\theta_{2}$ (avec $\theta_{2}>\theta_{1}$ ).
Quel est l'effet de l'augmentation de la température sur~:
a. la vitesse volumique de réaction~?
b. la valeur de l'avancement final~?

\subsection{Transformation acide-base}
À partir de la solution précédente $\mathrm{(S)}$, on prépare une solution
aqueuse $\mathrm{(S_{A})}$ de concentration molaire $C_{A}$. Le taux
d'avancement final de la réaction entre l'acide méthanoïque et l'eau dans ce cas
vaut $\tau=0,21$.

\begin{donnees}
    \begin{itemize}
        \item $\mathrm{pK_{A}}(\ce{HCO2H_{(aq)}}/\ce{HCO^-_{2(aq)}})=3,8$
    \end{itemize}
\end{donnees}
\begin{enumerate}
    \item Écrire l'équation chimique de la réaction de l'acide méthanoïque avec
          l'eau.
    \item Montrer que $C_{A}=\frac{1-\tau}{\tau^{2}}\cdot10^{-\mathrm{pK_{A}}}$.
          Calculer la valeur de $C_{A}$.
    \item Déterminer la valeur du $\mathrm{pH}$ de la solution
          $\mathrm{(S_{A})}$.
    \item En utilisant le dispositif schématisé sur la
          Fig.~\ref{fig:2bac_svt_national_2023_rattrapage_chim_1}, on dose le
          volume $V_{A}=\qty{30}{\mL}$ de la solution $\mathrm{(S_{A})}$ par une
          solution aqueuse $\mathrm{(S_{B})}$ d'hydroxyde de sodium
          $\ce{Na^+_{(aq)} + HO^-_{(aq)}}$de concentration molaire
          $C_{B}=\qty{1e-2}{\mol\per\L}$. La courbe de la
          Fig.~\ref{fig:2bac_svt_national_2023_rattrapage_chim_2}
          représente la variation du $\mathrm{pH}$ du mélange en fonction du
          volume $V_{B}$ de la solution $\mathrm{(S_{B})}$ versé au cours du
          dosage.

          \begin{minipage}{0.49\linewidth}
              \begin{figure}[H]
                  \centering
                  \includegraphics[width=0.8\textwidth]{2026-03-06_115446-}
                  \caption{}
                  \label{fig:2bac_svt_national_2023_rattrapage_chim_1}
              \end{figure}
          \end{minipage}
          \hfill
          \begin{minipage}{0.49\linewidth}
              \begin{figure}[H]
                  \centering
                  \includegraphics[width=0.9\textwidth]{2026-03-06_115432-}
                  \caption{}
                  \label{fig:2bac_svt_national_2023_rattrapage_chim_2}
              \end{figure}
          \end{minipage}
          
          \begin{enumerate}
              \item Nommer les éléments du dispositif numérotés 1, 2 et 3.
              \item Déterminer les coordonnées du point d'équivalence.
              \item Retrouver la valeur de $C_{A}$.
              \item Préciser, parmi les indicateurs colorés ci-dessous,
                    l'indicateur convenable pour repérer cette équivalence.
          
                    \begin{minipage}{\linewidth}
                        \begin{table}[H]
                            \centering
                            \begin{tabularx}{\textwidth}{|p{4cm}|C|C|C|}
                                \hline
                                \textbf{Indicateur coloré} &
                                \textbf{Teinte acide} &
                                \textbf{Zone de virage} &
                                \textbf{Teinte basique} \\
                                \hline
                                Rouge de méthyle & Rouge & $4,4-6,2$ & Jaune \\
                                \hline
                                Rouge de crésol & Jaune & $7,2-8,8$ & Rouge \\
                                \hline
                                Alizarine & Rouge & $11,0-12,4$ & Violette \\
                                \hline
                            \end{tabularx}
                        \end{table}
                    \end{minipage}
          \end{enumerate}
\end{enumerate}



\subsection{Transformation d'oxydo-réduction}
Pour étudier la réaction entre le cuivre et le dibrome en solution, on place
dans un bécher~:
\begin{itemize}
    \item le volume $V_{1}=\qty{50}{\mL}$ d'une solution de bromure de
          cuivre(II) $\ce{Cu^2+_{(aq)} + 2Br^-_{(aq)}}$ de concentration molaire
          ${C_{1}=\qty{4e-2}{\mol\per\L}}$~;
    \item le volume $V_{2}=\qty{50}{\mL}$ d'une solution de dibrome
          $\ce{Br2_{(aq)}}$ de concentration molaire
          ${C_{2}=\qty{4e-2}{\mol\per\L}}$~;
    \item la poudre de cuivre.
\end{itemize}

\begin{donnees}
    \begin{itemize}
        \item Constante d'équilibre associée à l'équation
              $\ce{Br2_{(aq)} + Cu_{(s)} <=>[(1)][(2)]
              2Br^-_{(aq)} + Cu^2+_{(aq)}}$~: $\mathrm{K}=1,2.10^{25}$
        \item Le cuivre est en excès.
    \end{itemize}
\end{donnees}
\begin{enumerate}
    \item Identifier les couples (ox/red) intervenants dans la réaction.
    \item Calculer la valeur du quotient de réaction $Q_{r,i}$ à l'état initial
          du système chimique.
    \item Préciser, en justifiant, le sens d'évolution du système chimique.
\end{enumerate}




